\uuid{Rhhp}
\titre{Polynômes/Système linéaire}
\auteur{hanani ; mzari}
\organisation{exo7}
\licence{CC BY-NC-SA 4.0}
\datecreate{2020-01-20}
\niveau{L1}
\module{Algèbre}
\chapitre{Déterminant, système linéaire}
\sousChapitre{Système linéaire, rang}
\difficulte{3}
\theme{Systèmes d'équations linéaires, Polynômes, Système linéaire}
\exo7qcmid{lille-2.systemes.023}

\contenu{
\begin{question}
\qclassification{200.03}
\qtags{motcle=Polynômes/Système linéaire}
Soit $P$ un polynôme à coefficients réels de degré $\le 3$ vérifiant les conditions :
$$P(1)=1,\quad P(0)=1,\quad P(-1)=-1\quad \mbox{et}\quad P'(1)= 3.$$
Quelles sont les assertions vraies ?
\begin{answers}
\bad{Un tel polynôme $P$ n'existe pas.}
\bad{Il existe une infinité de polynômes $P$ vérifiant ces conditions.}
\good{Il existe un unique polynôme $P$  vérifiant ces conditions.}
\bad{Si $P$ est un polynôme qui vérifie ces conditions, alors $P(2)=2$.}
\end{answers}
\begin{explanations} On pose : $P(X)=aX^3+bX^2+cX+d$, où $a,b,c$ et $d$ sont des réels à déterminer. En résolvant le système :
$$(\mathtt{S})
\left\{\begin{array}{rcc}
P(1)&=&1\\
P(0)&=&1\\
P(-1)&=&-1\\
P'(1)&=&3 \\
\end{array}\right.
\Leftrightarrow
\left\{\begin{array}{rcc}
a+b+c+d&=&1\\
d&=&1\\
-a+b-c+d&=&-1\\
3a+2b+c&=&3. \\
\end{array}\right. $$
On obtient : $P(X)=2X^3-X^2-X+1$. Par conséquent il existe un unique polynôme vérifiant les conditions ci-dessus et $P(2)=11$.
\end{explanations}
\end{question}
}
